\subsection{Full-Budget Benchmark Protocol}
The final benchmark was run with $10$ random seeds, namely $7$ through $16$, with one CUDA job per seed. The benchmark consisted of two blocks. The first block covered AWGN, Rayleigh, and SSPA. The second block covered the optical-fiber channel OptFib through a separate benchmark path. In the final suite, we retained only the following methods: drifting in direct-output mode, drifting in residual mode, DDPM, DDIM, and WGAN. We did not include the later \texttt{GAN\_FA} baseline or the later standalone GAN baseline in this final comparison.

For AWGN, Rayleigh, and SSPA we used channel-specific presets from the direct-output benchmark runner. AWGN and Rayleigh used $n=7$, while SSPA used $n=8$. The corresponding $(E_b/N_0, R)$ pairs were $(5\,\mathrm{dB}, 4/7)$ for AWGN, $(12\,\mathrm{dB}, 4/7)$ for Rayleigh, and $(8\,\mathrm{dB}, 6/8)$ for SSPA, which yielded noise standard deviations $0.5260$, $0.2350$, and $0.3251$, respectively. AWGN and Rayleigh were trained with dataset size $10^7$, batch size $5000$, and $30$ epochs for drifting, diffusion, and WGAN. SSPA used dataset size $10^7$, batch size $4096$, and $160$ epochs for all three method families. Although the benchmark preset defines an evaluation size of $10^7$, the full suite used an evaluation size of $10^6$ per method for tractable multi-seed execution. All results on AWGN, Rayleigh, and SSPA used $128$ SWD projections.

The drifting model used the same architecture on all channels: a conditional MLP with latent dimension $16$, hidden width $128$, two hidden layers, and \texttt{SiLU} activations, trained with Adam at learning rate $10^{-3}$. The drifting loss used drift scale $1.0$, median-heuristic bandwidth selection, minimum bandwidth $0.2$, maximum drift norm $2.0$, and repulsive weight $1.0$. We trained both parameterizations: direct drifting with target $y$ and residual drifting with target $e=y-x$. On AWGN, Rayleigh, and SSPA, direct drifting was evaluated in output space and residual drifting in residual space.

The diffusion baseline used direct-output modeling, i.e.\ $p(y\mid x)$, with $v$-prediction, a cosine-ZF variance schedule, $100$ diffusion steps, EMA decay $0.9$, and the conditional MLP backbone from our benchmark runner. The hidden width was $110$ for AWGN, $128$ for Rayleigh, and $110$ for SSPA. AWGN and Rayleigh used the two-stage learning-rate schedule from the reference setup, namely $10$ epochs at $10^{-3}$ followed by $20$ epochs at $10^{-4}$. SSPA used a constant diffusion learning rate of $10^{-4}$ over $160$ epochs. We reported DDPM and DDIM with $100$, $50$, $20$, and $10$ denoising steps. On AWGN, Rayleigh, and SSPA, diffusion was evaluated in output space.

WGAN used the conditional Wasserstein formulation with RMSprop at learning rates $10^{-4}$ for both generator and critic, five critic steps per generator step, and critic weight clipping at $0.01$. The hidden width was $128$ for AWGN and $256$ for Rayleigh and SSPA. The conditioning variable was normalized internally before training. In the current benchmark implementation, WGAN reports SWD in residual space.

OptFib was benchmarked separately with dataset size $120{,}000$, batch size $512$, $60$ training epochs, evaluation size $100{,}000$, and one DDIM setting, namely DDIM-$100$. The optical-fiber channel used the parameters $L=5000$, $\gamma=1.27$, $K_{\text{step}}=50$, and $P_n=-21.3$ dBm, with \texttt{use\_noise\_std=False}. Consequently, the effective channel noise was fixed by the optical-noise model and not by the nominal \texttt{noise\_std} field stored in the training configs. For drifting and diffusion in the OptFib path, the effective signal dimension was $n=2$. OptFib drifting again used both direct and residual parameterizations with the same drifting hyperparameters as above. OptFib diffusion used the residual formulation, $\epsilon$-prediction, the standard cosine schedule, $100$ diffusion steps, hidden width $128$, learning rate $10^{-3}$, EMA decay $0.995$, and $256$ SWD projections. OptFib WGAN used hidden width $128$, RMSprop learning rates $10^{-4}$, five critic steps, clipping at $0.01$, and $256$ SWD projections.

All seeds were initialized through a deterministic seeding routine. For AWGN, Rayleigh, and SSPA, the SWD random-projection seed was set equal to the current run seed. For the OptFib block, the optional-baseline runner used the fixed metric seed $12345$. Final benchmark tables report the mean and standard deviation over the ten seeds.
